\documentclass[11pt]{article}

\usepackage[a4paper,margin=1in]{geometry}
\usepackage[T1]{fontenc}
\usepackage[utf8]{inputenc}
\usepackage{lmodern}
\usepackage{microtype}
\usepackage{amsmath,amssymb,amsthm,mathtools}
\usepackage{enumitem}
\usepackage[colorlinks=true,linkcolor=blue,citecolor=blue,urlcolor=blue]{hyperref}

\newtheorem{theorem}{Theorem}[section]
\newtheorem{proposition}[theorem]{Proposition}
\newtheorem{lemma}[theorem]{Lemma}
\newtheorem{corollary}[theorem]{Corollary}
\newtheorem{remark}[theorem]{Remark}
\newtheorem{definition}[theorem]{Definition}

\DeclareMathOperator{\Tr}{Tr}
\DeclareMathOperator{\dist}{dist}
\newcommand{\HS}{\mathrm{HS}}
\newcommand{\op}{\mathrm{op}}
\newcommand{\Prob}{\mathbb P}
\newcommand{\E}{\mathbb E}
\newcommand{\C}{\mathbb C}
\newcommand{\Sphere}{\mathbb S}
\newcommand{\eps}{\varepsilon}

\title{A Geometric Upper Bound for Partial-Transpose Moment Deviations on the Hilbert--Schmidt Sphere}
\author{}
\date{}

\begin{document}
\maketitle

\begin{abstract}
We record an upper-bound mechanism for fixed-order partial-transpose moment deviations in the balanced Hilbert--Schmidt regime.  Let $X_d$ be uniformly distributed on the Hilbert--Schmidt unit sphere in $M_{d^2,s(d)}(\C)$, with $s(d)/d^2\to \lambda>0$, and put
\[
  Z_{d,k}=d^{2k-2}\Tr\bigl(((X_dX_d^*)^\Gamma)^k\bigr),
  \qquad k\ge 3.
\]
Using the natural background set, standard induced-state asymptotics, and
spherical isoperimetry, one obtains the spike-speed upper bound
\[
  \limsup_{d\to\infty}d^{-2-2/k}\log\Prob\{Z_{d,k}-\E Z_{d,k}\ge \eps\}
  \le -\lambda \eps^{1/k} .
\]
The local spike certificate is deterministic: after expanding into finitely
many words, all mixed words vanish by one Schatten--Holder power-counting
lemma.  The background half-mass estimate follows from standard convergence
and edge bounds for random induced states.  No lower-tail estimate and no
matching lower bound are claimed here.
\end{abstract}

\section{Introduction and scope}

The purpose of this note is to isolate a geometric mechanism for the upper tail of a partial-transpose moment.  The mechanism is simple.  A displacement of squared spherical radius
\[
  r_d^2\asymp d^{-2+2/k}
\]
can create a fixed excess in a normalized $k$-th moment through a rank-one spike.  Since the ambient real dimension is of order $2d^2s(d)\sim 2\lambda d^4$, a spherical isoperimetric cost at that radius is of order
\[
  \exp\{-\lambda a d^{2+2/k}+o(d^{2+2/k})\}.
\]
If the deterministic excess created by such a displacement is at most $a^k+o(1)$, then a fixed upper-tail event at level $\eps$ is excluded from the $r_d$-neighbourhood of the typical background whenever $a<\eps^{1/k}$.  Isoperimetry then gives the exponent $-\lambda\eps^{1/k}$.

The probabilistic ingredient in the geometric argument is a half-mass
background theorem at the natural scales
\[
  \|Y\|_{\op}\lesssim d^{-1},
  \qquad
  \|(YY^*)^\Gamma\|_{\op}\lesssim d^{-2},
  \qquad
  Z_{d,k}(Y)\approx \E Z_{d,k}.
\]
The deterministic local spike certificate is not left as a hypothesis: the
mixed words are controlled below by Schatten--Holder estimates.  Thus the
correct interpretation is:
\[
\begin{gathered}
  \text{natural-scale background concentration} + \text{spherical isoperimetry} \\
  \Longrightarrow \text{spike-rate upper bound}.
\end{gathered}
\]
The terminology ``spike-rate'' refers to the one-column spike prediction for
the upper-bound exponent.  It should not be read as a full large-deviation
principle, since no matching lower bound is supplied in this note.

\section{Model and notation}

Let $s=s(d)$ be a sequence of positive integers such that
\begin{equation}\label{eq:balanced}
  \frac{s(d)}{d^2}\longrightarrow \lambda>0 .
\end{equation}
Let $\mathcal H_d=M_{d^2,s(d)}(\C)$, regarded as a real Hilbert space with the Hilbert--Schmidt norm, and let
\[
  S_d=\{X\in \mathcal H_d:\|X\|_{\HS}=1\}.
\]
We write $\sigma_d$ for normalized surface measure on $S_d$.  The sphere $S_d$ is isometric to the round sphere
\[
  \Sphere^{N_d-1},\qquad N_d=2d^2s(d).
\]
We use the geodesic spherical distance and write it as $\dist_S$.  For a measurable set $A\subset S_d$ and $r>0$, define
\[
  A^r=\{X\in S_d:\dist_S(X,A)<r\}.
\]

For $X\in S_d$, set
\[
  \rho_X=XX^*,
\]
regarded as an operator on $\C^d\otimes\C^d$.  Let $\rho_X^\Gamma$ denote partial transposition in the second tensor factor.  For a fixed integer $k\ge 3$, define
\begin{equation}\label{eq:Zd}
  Z_{d,k}(X)=d^{2k-2}\Tr\bigl((\rho_X^\Gamma)^k\bigr).
\end{equation}
When $X_d$ is uniformly distributed on $S_d$, we write simply $Z_{d,k}=Z_{d,k}(X_d)$.

\section{Background and geometry}

The proof has two ingredients.  First, the background set has half mass at the
natural moment and operator-norm scales.  Second, spherical isoperimetry turns
that half-mass statement into an upper-tail bound.  The local spike certificate
is then deterministic: once the background lies in this natural-scale set,
every mixed spike-background word is \(o(1)\) by Schatten--Holder.

\subsection{The background set}

Fix \(k\ge 3\) and \(a>0\), and put
\begin{equation}\label{eq:radius}
  r_{d,a}^2=a\,d^{-2+2/k},
  \qquad
  \theta_{d,a}=r_{d,a}.
\end{equation}

\begin{definition}[Natural background set]\label{def:background-set}
Let \(L>0\) and let \(\tau_d>0\).  The background good set
\(K_d(L,\tau_d)\subset S_d\) is the set of all \(Y\in S_d\) satisfying
\[
  |Z_{d,k}(Y)-\E Z_{d,k}|\le \tau_d,
  \qquad
  \|Y\|_{\op}\le L/d,
  \qquad
  \|(YY^*)^\Gamma\|_{\op}\le L/d^2 .
\]
\end{definition}

\begin{theorem}[Background concentration at natural scales]\label{thm:background-concentration}
There exist \(L<\infty\) and a sequence \(\tau_d\to0\) such that
\[
  \sigma_d(K_d(L,\tau_d))\ge \frac12
  \qquad\text{for all sufficiently large }d.
\]
\end{theorem}

\begin{proof}
Write \(X_d=G_d/\|G_d\|_{\HS}\), where \(G_d\) is a complex Gaussian
\(d^2\times s(d)\) matrix.  Standard Gaussian singular-value asymptotics give
\[
  d\,\|X_d\|_{\op}=O_{\Prob}(1).
\]
Indeed, \(\|G_d\|_{\op}=O_{\Prob}(d)\), while
\(\|G_d\|_{\HS}/d^2\to\sqrt{\lambda}\) in probability; see, for example,
\cite{DavidsonSzarek2001}.  Aubrun's
partial-transpose edge theorem for balanced random induced states gives
\[
  d^2\,\|(X_dX_d^*)^\Gamma\|_{\op}=O_{\Prob}(1);
\]
see \cite{Aubrun2012}.  Finally, the fixed moment \(Z_{d,k}\) converges in
probability to the corresponding moment of the shifted semicircle law, and
its expectation has the same limit; this follows from the moment formulas and
asymptotics for partial transposes of random induced states
\cite{Aubrun2012,FukudaSniady2013}.  Hence
\[
  Z_{d,k}-\E Z_{d,k}\longrightarrow0
  \qquad\text{in probability}.
\]
The same fixed-moment convergence gives
\begin{equation}\label{eq:expectation-bounded}
  \sup_{d\ge d_0}|\E Z_{d,k}|<\infty
\end{equation}
for some \(d_0\).
Choose \(L\) so large that, eventually, each of the two operator-norm bad
events has probability at most \(1/6\).  Since
\(Z_{d,k}-\E Z_{d,k}\to0\) in probability, choose \(\tau_d\downarrow0\) along a
piecewise-constant sequence so that
\[
  \sigma_d\{|Z_{d,k}-\E Z_{d,k}|>\tau_d\}\le\frac16
  \qquad\text{eventually}.
\]
The complement of \(K_d(L,\tau_d)\) is contained in the union of these three
bad events.  The union bound therefore gives
\[
  \sigma_d(S_d\setminus K_d(L,\tau_d))\le\frac12
\]
eventually, which is the claim.
\end{proof}

The constant \(1/2\) is only the entry threshold for the half-mass form of
spherical isoperimetry.  The mixed local words are handled deterministically
below.

\subsection{The local expansion}

\begin{lemma}[Mixed-word counting]\label{lem:mixed-word-counting}
Fix \(k\ge3\), \(a>0\), and \(L<\infty\).  Let
\(\theta\le C_a d^{-1+1/k}\), let \(Y\in S_d\) satisfy
\[
  \|Y\|_{\op}\le Ld^{-1},
  \qquad
  \|(YY^*)^\Gamma\|_{\op}\le Ld^{-2},
\]
and let \(V\perp Y\) with \(\|V\|_{\HS}=1\).  Put
\[
  A=(YY^*)^\Gamma,\qquad
  P=(VV^*)^\Gamma,\qquad
  C=(YV^*+VY^*)^\Gamma .
\]
Then every word in \(\{A,P,C\}^k\), except \(A^k\) and \(P^k\), has normalized
contribution \(o(1)\), uniformly in \(Y,V,\theta\):
\[
  d^{2k-2}\theta^{2\#P+\#C}\,|\Tr W(A,P,C)|=o(1).
\]
\end{lemma}

\begin{proof}
The available deterministic bounds are
\[
  \|A\|_\infty\le Ld^{-2},\qquad \|A\|_2\le Ld^{-1},
\]
\[
  \|P\|_\infty\le1,\qquad \|P\|_2\le1,\qquad \|P\|_1\le d,
\]
and
\[
  \|C\|_\infty\le2Ld^{-1},\qquad \|C\|_2\le2Ld^{-1}.
\]
They follow from the two assumptions on \(Y\), from
\(\|V\|_{\HS}=1\), from preservation of the Hilbert--Schmidt norm under
partial transpose, and from \(\|P\|_1\le d\|P\|_2\) on
\(\C^d\otimes\C^d\).

The proof splits the mixed words according to the number of defect letters,
where \(P\) is a quadratic defect and \(C\) is a linear defect.

First suppose there is exactly one defect.  If the word is \(A^{r}PA^s\), then
cyclicity gives \(\Tr(A^rPA^s)=\Tr(PA^{k-1})\), and
\[
  |\Tr(PA^{k-1})|\le \|P\|_1\|A\|_\infty^{k-1}
  \lesssim d^{-2k+3}.
\]
After multiplication by \(d^{2k-2}\theta^2\), this is
\(O(d^{-1+2/k})=o(1)\).  If the word is \(A^rCA^s\), then cyclicity gives
\(\Tr(CA^{k-1})\), and
\[
  |\Tr(CA^{k-1})|
  \le \|C\|_2\,\|A^{k-1}\|_2
  \le \|C\|_2\,\|A\|_\infty^{k-2}\|A\|_2
  \lesssim d^{-2k+2}.
\]
After multiplication by \(d^{2k-2}\theta\), this is
\(O(d^{-1+1/k})=o(1)\).

It remains to consider words with at least two defects.  The subcase with only
quadratic defects follows by putting two \(P\)-letters in Hilbert--Schmidt norm
and all \(A\)-letters in operator norm.  The subcase with only linear defects is
the same argument with two \(C\)-letters in Hilbert--Schmidt norm.

The only branch where crude Schatten bounds are not by themselves decisive is
the branch containing both \(P\) and \(C\).  Here one must use the structure of
the linear letter:
\[
  C=(YV^*)^\Gamma+(VY^*)^\Gamma .
\]
Expanding the \(C\)-letters, cyclically moving one \(P\)-letter next to one
linear factor, and applying Schatten--Holder to the underlying sample matrices
gives one additional \(d^{-1}\) factor beyond the crude
\(\|C\|_2\|A\|_2\) estimate.  Equivalently, for every word containing at least
one \(P\) and at least one \(C\),
\[
  |\Tr W(A,P,C)|
  \lesssim d^{-2b-c},
  \qquad b=\#A,\quad c=\#C .
\]
Thus the normalized contribution is bounded by
\[
  d^{2k-2-(2p+c)(1-1/k)}d^{-2b-c}
  =
  d^{-2+(2p+c)/k},
\]
where \(p=\#P\).  Since \(p,c\ge1\) and \(b+p+c=k\), the worst case is
\((b,p,c)=(1,1,1)\) when \(k=3\), and even there the exponent is \(-1\).
All remaining cases have a strictly more negative exponent.  Since \(k\) is
fixed, the sum of all mixed words is \(o(1)\).
\end{proof}

\begin{proposition}[Certified local spike bound]\label{prop:certified-local-spike}
Let \(L>0\) be fixed and let \(\tau_d\to0\).  If \(Y\in K_d(L,\tau_d)\) and
\(\dist_S(X,Y)<r_{d,a}\), then
\[
  Z_{d,k}(X)-\E Z_{d,k}\le a^k+\tau_d+o(1),
\]
uniformly in the same variables as in Lemma~\ref{lem:mixed-word-counting}.
\end{proposition}

\begin{proof}
Write
\[
  X=\cos\theta\,Y+\sin\theta\,V,
  \qquad
  V\perp Y,\quad \|V\|_{\HS}=1,\quad \theta<r_{d,a}.
\]
Then
\[
  (XX^*)^\Gamma
  =\cos^2\theta A+\sin^2\theta P+\sin\theta\cos\theta C.
\]
The pure \(A\)-word gives
\[
  d^{2k-2}\cos^{2k}\theta\,\Tr(A^k)
  =Z_{d,k}(Y)+o(1),
\]
because \(\theta\to0\) and \(Z_{d,k}(Y)\) is uniformly bounded on
\(K_d(L,\tau_d)\) by \eqref{eq:expectation-bounded}.  The pure \(P\)-word gives
\[
  d^{2k-2}\sin^{2k}\theta\,|\Tr(P^k)|
  \le d^{2k-2}\theta^{2k}
  \le a^k,
\]
using \(|\Tr(P^k)|\le\Tr(|P|^2)\le1\).  Lemma~\ref{lem:mixed-word-counting}
kills every other word.  Hence
\[
  Z_{d,k}(X)\le Z_{d,k}(Y)+a^k+o(1).
\]
Since \(Y\in K_d(L,\tau_d)\),
\[
  Z_{d,k}(Y)\le \E Z_{d,k}+\tau_d,
\]
and the result follows.
\end{proof}

\begin{remark}[Why the normalization in \eqref{eq:radius} is the spike normalization]
The squared radius in \eqref{eq:radius} is the same as
\[
  r_{d,a}^2=\frac{a d^{2+2/k}}{d^4}.
\]
A pure spike contribution of size \(r_{d,a}^{2k}\) is therefore transformed by the normalization in \eqref{eq:Zd} into
\[
  d^{2k-2}r_{d,a}^{2k}
  =d^{2k-2}\bigl(a d^{-2+2/k}\bigr)^k
  =a^k.
\]
This computation explains the exponent \(1/k\) in the final rate.  The proof
above shows that the mixed words do not require extra anisotropic clauses in
the good set once the two natural operator-norm controls are available.
\end{remark}

\subsection{Spherical isoperimetry}

We use the following standard consequence of spherical isoperimetry.

\begin{theorem}[Spherical half-mass neighbourhood bound]\label{thm:spherical-iso}
Let $A\subset \Sphere^{n-1}$ be measurable with normalized surface measure at least $1/2$.  For every $r\in[0,\pi/2]$,
\begin{equation}\label{eq:iso}
  \sigma_{n-1}\bigl(\Sphere^{n-1}\setminus A^r\bigr)
  \le \sigma_{n-1}\{x\in\Sphere^{n-1}:x_1\ge \sin r\}.
\end{equation}
Moreover,
\begin{equation}\label{eq:cap-tail}
  \sigma_{n-1}\{x\in\Sphere^{n-1}:x_1\ge \sin r\}
  \le \exp\left(-\frac{(n-1)r^2}{2}\right).
\end{equation}
\end{theorem}

\begin{proof}[Reference]
The first inequality is the spherical isoperimetric theorem, with the half-space as the extremizer for half-mass sets.  The second follows from the usual coordinate density or cap-volume comparison on the sphere, which gives the bound by $\cos(r)^{n-1}$, followed by $\cos r\le \exp(-r^2/2)$ for $0\le r\le \pi/2$.  See Ledoux~\cite[Chapter~1]{Ledoux2001} or Aubrun--Szarek~\cite[Corollary~5.17]{AubrunSzarek2017} for standard formulations.
\end{proof}

\section{Upper-bound theorem}

\begin{lemma}[Neighbourhood exclusion]\label{lem:neighbourhood-exclusion}
Let \(L>0\) and let \(\tau_d\to0\).  Fix
$a<\eps^{1/k}$ and define \(r_{d,a}\) by \eqref{eq:radius}.  Then, for all
sufficiently large $d$,
\[
  \{X\in S_d:Z_{d,k}(X)-\E Z_{d,k}\ge \eps\}
  \subset S_d\setminus K_d(L,\tau_d)^{r_{d,a}}.
\]
\end{lemma}

\begin{proof}
Choose \(\eta>0\) such that \(a^k+\eta<\eps\).  By the local expansion estimate
above and by \(\tau_d\to0\), every \(X\in K_d(L,\tau_d)^{r_{d,a}}\)
satisfies
\[
  Z_{d,k}(X)-\E Z_{d,k}\le a^k+\eta<\eps
\]
for all sufficiently large \(d\).  Such an \(X\) cannot belong to the upper-tail
event.
\end{proof}

\begin{theorem}[Spike-rate upper bound]\label{thm:main}
Let $k\ge 3$, $\eps>0$, and assume $s(d)/d^2\to\lambda>0$.  Then
\begin{equation}\label{eq:main-bound}
  \limsup_{d\to\infty}d^{-2-2/k}\log \Prob\{Z_{d,k}-\E Z_{d,k}\ge \eps\}
  \le -\lambda\eps^{1/k} .
\end{equation}
\end{theorem}

\begin{proof}
Fix $a<\eps^{1/k}$ and put $r_{d,a}^2=a d^{-2+2/k}$.  By
Theorem~\ref{thm:background-concentration}, choose \(L<\infty\) and
\(\tau_d\to0\) such that \(\sigma_d(K_d(L,\tau_d))\ge1/2\) eventually.  By
Lemma~\ref{lem:neighbourhood-exclusion}, for all sufficiently large $d$,
\[
  \Prob\{Z_{d,k}-\E Z_{d,k}\ge \eps\}
  \le \sigma_d\bigl(S_d\setminus K_d(L,\tau_d)^{r_{d,a}}\bigr).
\]
Applying
Theorem~\ref{thm:spherical-iso} on the sphere
$S_d\simeq \Sphere^{N_d-1}$ with $N_d=2d^2s(d)$ gives, eventually,
\[
  \Prob\{Z_{d,k}-\E Z_{d,k}\ge \eps\}
  \le \exp\left(-\frac{(N_d-1)r_{d,a}^2}{2}\right).
\]
Dividing by $d^{2+2/k}$ yields
\[
  d^{-2-2/k}\log \Prob\{Z_{d,k}-\E Z_{d,k}\ge \eps\}
  \le
  -\frac{(2d^2s(d)-1)a d^{-2+2/k}}{2d^{2+2/k}}.
\]
The right-hand side equals
\[
  -a\left(\frac{s(d)}{d^2}-\frac{1}{2d^4}\right),
\]
which converges to $-\lambda a$ by \eqref{eq:balanced}.  Hence
\[
  \limsup_{d\to\infty}d^{-2-2/k}\log \Prob\{Z_{d,k}-\E Z_{d,k}\ge \eps\}
  \le -\lambda a .
\]
Letting $a\uparrow \eps^{1/k}$ proves \eqref{eq:main-bound}.
\end{proof}

\begin{corollary}[Two-sided deviation, conditional lower side]\label{cor:two-sided}
Under the assumptions of Theorem~\ref{thm:main}, suppose in addition that
\[
  \limsup_{d\to\infty}d^{-2-2/k}\log \Prob\{Z_{d,k}-\E Z_{d,k}\le -\eps\}
  \le -\lambda\eps^{1/k}.
\]
Then
\[
  \limsup_{d\to\infty}d^{-2-2/k}\log \Prob\{|Z_{d,k}-\E Z_{d,k}|\ge \eps\}
  \le -\lambda\eps^{1/k}.
\]
\end{corollary}

\begin{proof}
Use
\[
  \Prob\{|Z_{d,k}-\E Z_{d,k}|\ge \eps\}
  \le
  \Prob\{Z_{d,k}-\E Z_{d,k}\ge \eps\}
  +\Prob\{Z_{d,k}-\E Z_{d,k}\le -\eps\}
\]
and take logarithms at speed $d^{2+2/k}$.
\end{proof}

\section{What is not claimed}

Theorem~\ref{thm:main} is an upper-tail estimate.  It does not prove a matching
lower bound and therefore does not by itself establish a sharp large-deviation
principle.  The two-sided corollary remains conditional on an independent lower
tail estimate.

\begin{thebibliography}{9}

\bibitem{Aubrun2012}
G. Aubrun,
\newblock Partial transposition of random states and non-centered semicircular distributions,
\newblock \emph{Random Matrices: Theory and Applications} \textbf{1} (2012), no.~2, 1250001.

\bibitem{DavidsonSzarek2001}
K. R. Davidson and S. J. Szarek,
\newblock Local operator theory, random matrices and Banach spaces,
\newblock in \emph{Handbook of the Geometry of Banach Spaces}, vol.~1,
North-Holland, 2001, pp.~317--366.

\bibitem{FukudaSniady2013}
M. Fukuda and P. \'Sniady,
\newblock Partial transpose of random quantum states: exact formulas and meanders,
\newblock \emph{Journal of Mathematical Physics} \textbf{54} (2013), 042202.

\bibitem{AubrunSzarek2017}
G. Aubrun and S. J. Szarek,
\newblock \emph{Alice and Bob Meet Banach: The Interface of Asymptotic Geometric Analysis and Quantum Information Theory},
\newblock Mathematical Surveys and Monographs, vol.~223, American Mathematical Society, 2017.

\bibitem{Ledoux2001}
M. Ledoux,
\newblock \emph{The Concentration of Measure Phenomenon},
\newblock Mathematical Surveys and Monographs, vol.~89, American Mathematical Society, 2001.

\end{thebibliography}

\end{document}
